% sec-04-results.tex - Results (full)
\section{Results}

\subsection{A Single Prompt Produced Four Stage Outputs}

Executing the single master prompt produced four stage outputs in
\texttt{trial-documents}: 24 colored Mermaid figures, a draft scaffold, a full
paper, and a final paper. Each LaTeX stage ships \texttt{main.tex},
\texttt{paperstyle.sty}, \texttt{references.bib}, a \texttt{sections/} directory of
eight section files, and an Overleaf zip, and every distinguishable file was a
separate commit pushed in real time. Table 2 summarizes the stage outputs.

\vspace{0.2em}
\noindent\begin{tabularx}{\textwidth}{L{2.4cm} L{3.4cm} Y L{2.0cm}}
\toprule
\textbf{Stage} & \textbf{Directory} & \textbf{Artifacts} & \textbf{Commits} \\
\midrule
1 Mermaid & \texttt{mermaid/} & 24 colored Mermaid figure files, README, output-mermaid.md & 26 \\
2 Draft & \texttt{draft-paper/} & main.tex, paperstyle.sty, references.bib, 8 section scaffolds, README, zip & 10+ \\
3 Full & \texttt{full-paper/} & main.tex, style, bib, 8 full sections with 24 TikZ figures and tables, README, zip & 10+ \\
4 Final & \texttt{final-paper/} & polished main, style, bib, 8 sections, README, zip; root repository updates & 10+ \\
\bottomrule
\end{tabularx}
\figcaption{Table 2. The four stage outputs produced by the single prompt, each
with one commit per distinguishable file pushed in real time.}

\subsection{The Six Acceleration Targets Addressed}

The six document targets where faster, validated authoring yields the greatest
schedule value were each addressed by the method, with the binding clock identified
for each (Figure 16, Table 3). Two targets are detailed further: the initial IND
and IRB package, the single highest-value target (Figure 17), and the complete
clinical-hold response, whose 30-day FDA review begins only when a complete response
to every deficiency is received (Figure 18).

\begin{mermaidfig}
\node[mmin] (in) at (0,-3.75) {Faster, validated\\authoring};
\node[mmgoal] (t1) at (5.0,0) {1. IND + IRB package};
\node[mmgoal] (t2) at (5.0,-1.5) {2. Amendments + consent};
\node[mmaccent] (t3) at (5.0,-3.0) {3. Cohort-review pkg};
\node[mmgoal] (t4) at (5.0,-4.5) {4. Clinical-hold response};
\node[mmaccent] (t5) at (5.0,-6.0) {5. EOP2 + Phase 3 protocol};
\node[mmaccent] (t6) at (5.0,-7.5) {6. CSR + NDA/BLA modules};
\node[mmgoal] (out) at (10.4,-3.75) {Compressed\\admin/prep time\\(months to a year)};
\draw[mmedge] (in) to[out=60,in=180] (t1.west);
\draw[mmedge] (in) to[out=40,in=180] (t2.west);
\draw[mmedge] (in) to[out=20,in=180] (t3.west);
\draw[mmedge] (in) to[out=-20,in=180] (t4.west);
\draw[mmedge] (in) to[out=-40,in=180] (t5.west);
\draw[mmedge] (in) to[out=-60,in=180] (t6.west);
\draw[mmedge] (t1.east) to[out=0,in=120] (out.north);
\draw[mmedge] (t2.east) to[out=0,in=150] (out.north west);
\draw[mmedge] (t3.east)--(out.west);
\draw[mmedge] (t4.east)--(out.west);
\draw[mmedge] (t5.east) to[out=0,in=210] (out.south west);
\draw[mmedge] (t6.east) to[out=0,in=240] (out.south);
\end{mermaidfig}
\figcaption{Figure 16. The six greatest-acceleration document targets. One input,
faster validated authoring, feeds all six ranked targets, each contributing to a
common compression of administrative and preparation time.}

\vspace{0.2em}
\noindent\begin{tabularx}{\textwidth}{L{0.6cm} L{4.2cm} L{2.6cm} Y}
\toprule
\textbf{\#} & \textbf{Document target} & \textbf{Gate} & \textbf{Binding clock or constraint} \\
\midrule
1 & Initial IND and IRB package & Hard & 30-day FDA review; IRB meeting calendar; cannot generate missing tox or stability data \\
2 & Protocol amendments + synchronized consent/site updates & Hard & IRB approval; speed is lost if protocol, consent, and database are revised serially \\
3 & Cohort-review packages after safety data mature & Protocol-defined & The protocol-required DLT observation window \\
4 & Complete clinical-hold response & Hard & 30-day FDA review of a complete response to every deficiency \\
5 & Phase 2-to-3 briefing package and Phase 3 protocol & Decision & EOP2 meeting scheduling and review \\
6 & Pivotal CSR and NDA/BLA modules after database lock & Decision/filing & Database lock; query resolution; RTOR staging \\
\bottomrule
\end{tabularx}
\figcaption{Table 3. The six acceleration targets, their gate type, and the binding
clock that faster authoring cannot shorten (adapted from the document-types
sources).}

\begin{mermaidfig}
\node[mmgoal] (pkg) at (5.0,0) {Initial IND + IRB package\\(highest schedule value)};
\node[mmin] (c1) at (0,-2.0) {Clinical protocol};
\node[mmin] (c2) at (0,-3.4) {Investigator's Brochure};
\node[mmin] (c3) at (0,-4.8) {Nonclinical, CMC, stability};
\node[mmin] (c4) at (10.0,-2.0) {ICF, recruitment};
\node[mmin] (c5) at (10.0,-3.4) {Investigator forms};
\node[mmin] (c6) at (10.0,-4.8) {Safety information};
\node[mmproc] (clk) at (5.0,-3.4) {30-day FDA clock\\starts at submission};
\node[mmaccent] (acc) at (2.5,-6.2) {faster assembly\\starts clock + sites sooner};
\node[mmdec] (lim) at (7.5,-6.2) {cannot shorten\\30 days or make\\missing data};
\draw[mmedge] (c1) to[out=0,in=160] (pkg.west);
\draw[mmedge] (c2) to[out=0,in=180] (clk.west);
\draw[mmedge] (c3) to[out=0,in=200] (clk.west);
\draw[mmedge] (c4) to[out=180,in=20] (pkg.east);
\draw[mmedge] (c5) to[out=180,in=0] (clk.east);
\draw[mmedge] (c6) to[out=180,in=-20] (clk.east);
\draw[mmedge] (pkg)--(clk);
\draw[mmedge] (clk) to[out=-120,in=90] (acc.north);
\draw[mmedge] (clk) to[out=-60,in=90] (lim.north);
\end{mermaidfig}
\figcaption{Figure 17. Composition of the initial IND and IRB package. Faster,
internally consistent assembly starts the 30-day FDA clock and all sites sooner but
cannot shorten the fixed period or generate missing toxicology or stability data.}

\begin{mermaidfig}
\node[mmdec] (h) at (0,0) {FDA clinical hold:\\trial cannot resume};
\node[mmin] (r1) at (4.4,1.3) {Revised protocol};
\node[mmin] (r2) at (4.4,0) {New nonclinical / CMC};
\node[mmin] (r3) at (4.4,-1.3) {Risk + consent update};
\node[mmproc] (asm) at (8.4,0) {Complete response:\\every deficiency};
\node[mmproc] (rev) at (8.4,-2.0) {FDA review:\\30 calendar days};
\node[mmgoal] (g) at (8.4,-3.8) {Hold lifted;\\enrollment resumes};
\node[mmaccent] (acc) at (3.6,-3.0) {faster assembly\\moves restart forward};
\draw[mmedge] (h) to[out=30,in=180] (r1.west);
\draw[mmedge] (h) to[out=0,in=180] (r2.west);
\draw[mmedge] (h) to[out=-30,in=180] (r3.west);
\draw[mmedge] (r1) to[out=0,in=120] (asm.north west);
\draw[mmedge] (r2)--(asm);
\draw[mmedge] (r3) to[out=0,in=240] (asm.south west);
\draw[mmedge] (asm)--(rev); \draw[mmedge] (rev)--(g);
\draw[mmedged] (acc) to[out=0,in=180] (rev.west);
\end{mermaidfig}
\figcaption{Figure 18. The complete clinical-hold response. The 30-day FDA review
begins only when a complete response to every deficiency is received, so faster
assembly moves the potential restart date forward; an incomplete response does not
start the useful clock.}

\subsection{New Document Iterations in One to Four Days}

A new iteration of any large document completed in one to four days, against a
weeks-to-months baseline. The gain is confined to the administrative and
preparation bucket: the clinical and operational bucket (waiting for survival
endpoints to mature) and the fixed regulatory review bucket are unchanged
(Figure 19). Cumulatively, compressing only the writing bucket can save months, and
sometimes more than a year, across a program \cite{srcDocTypesGem}.

\begin{mermaidfig}
\node[mmdec] (b1) at (0,0) {Clinical / operational\\NOT compressible};
\node[mmaccent] (b2) at (5.0,0) {Administrative / prep\\COMPRESSIBLE};
\node[mmdec] (b3) at (10.0,0) {Regulatory review\\FIXED};
\node[mmctx] (b1d) at (0,-1.9) {recruitment, treatment,\\PFS/OS maturation};
\node[mmin] (b2d) at (5.0,-1.9) {data cleaning, analysis,\\document writing};
\node[mmctx] (b3d) at (10.0,-1.9) {30-day IND review,\\30-day hold review};
\node[mmproc] (llm) at (5.0,-3.7) {Repository LLM\\(acts on this bucket only)};
\node[mmgoal] (g) at (5.0,-5.4) {months to a year\\saved cumulatively};
\draw[mmedge] (b1)--(b1d); \draw[mmedge] (b2)--(b2d); \draw[mmedge] (b3)--(b3d);
\draw[mmedge] (llm)--(b2d);
\draw[mmedge] (llm)--(g);
\end{mermaidfig}
\figcaption{Figure 19. The three timeline buckets. Only the administrative and
preparation bucket compresses; the clinical/operational bucket and the fixed
regulatory review clocks do not. The repository LLM acts on the middle bucket only.}

\subsection{Effective Verifications: Grounding and Name-Matching}

Five checks make the build verifiable (Figure 20, Table 4): grounding via the 24
Mermaid figures; figures matching the surrounding context; real-time GitHub
visibility of each file; repository and directory name matching; and file-name
matching. The quality gates that prevent a document from generating a rework cycle
(amendments, regulator questions, IRB revisions, or an extra review cycle) are shown
in Figure 21.

\begin{mermaidfig}
\node[mmgoal] (g) at (9.4,-3.0) {Grounded,\\monitorable paper};
\node[mmdec] (k1) at (0,0) {1 Grounding};
\node[mmdec] (k2) at (0,-1.5) {2 Figures match context};
\node[mmdec] (k3) at (0,-3.0) {3 GitHub real-time};
\node[mmdec] (k4) at (0,-4.5) {4 Repo/dir names match};
\node[mmdec] (k5) at (0,-6.0) {5 File names match};
\node[mmin] (e1) at (4.7,0) {24 Mermaid figures};
\node[mmin] (e2) at (4.7,-1.5) {cited where discussed};
\node[mmin] (e3) at (4.7,-3.0) {one commit per file};
\node[mmin] (e4) at (4.7,-4.5) {trial-documents/...};
\node[mmin] (e5) at (4.7,-6.0) {sec-04-results.tex};
\draw[mmedge] (k1)--(e1); \draw[mmedge] (k2)--(e2); \draw[mmedge] (k3)--(e3);
\draw[mmedge] (k4)--(e4); \draw[mmedge] (k5)--(e5);
\draw[mmedge] (e1) to[out=0,in=150] (g.north west);
\draw[mmedge] (e2) to[out=0,in=170] (g.west);
\draw[mmedge] (e3)--(g.west);
\draw[mmedge] (e4) to[out=0,in=190] (g.west);
\draw[mmedge] (e5) to[out=0,in=210] (g.south west);
\end{mermaidfig}
\figcaption{Figure 20. The five name-matching and grounding verifications, each with
a concrete example, that together make the paper grounded and monitorable.}

\vspace{0.2em}
\noindent\begin{tabularx}{\textwidth}{L{0.6cm} L{4.6cm} Y}
\toprule
\textbf{\#} & \textbf{Verification check} & \textbf{Concrete example} \\
\midrule
1 & Grounding via Mermaid figures & The 24 figures in \texttt{mermaid/} carry the document counts and review clocks reused in the prose \\
2 & Figures match paper context & Each figure is cited in the sentence that discusses it (Figures 1 to 24) \\
3 & AI files viewable on GitHub in real time & One commit per file is pushed at generation, so each artifact appears on the branch as it is made \\
4 & Repository and directory names match & The paper refers to \texttt{trial-documents/mermaid} and \texttt{full-paper/sections}, which are the directories on disk \\
5 & File names match repo file names & A cited \texttt{sec-04-results.tex} is this file; \texttt{fig-04-acceleration-targets.md} is the figure file \\
\bottomrule
\end{tabularx}
\figcaption{Table 4. The five verification checks (OUTLINE theme 2) with a concrete
example for each.}

\begin{mermaidfig}
\node[mmin] (d) at (0,0) {Draft document};
\node[mmdec] (q1) at (3.3,0) {consistent?};
\node[mmdec] (q2) at (6.6,0) {validated data?};
\node[mmdec] (q3) at (9.9,0) {clickable DOI/URL,\\single dashes?};
\node[mmgoal] (g) at (9.9,-2.2) {Submission-ready\\(no extra cycle)};
\node[mmaccent] (rev) at (5.0,-2.2) {Revise};
\node[mmctx] (note) at (3.3,-3.8) {poor docs add amendments,\\regulator questions,\\IRB revisions, extra cycle};
\draw[mmedge] (d)--(q1); \draw[mmedge] (q1)--(q2); \draw[mmedge] (q2)--(q3);
\draw[mmedge] (q3)--(g);
\draw[mmedged] (q1) to[out=-90,in=60] (rev.north east);
\draw[mmedged] (q2) to[out=-90,in=90] (rev.north);
\draw[mmedged] (rev) to[out=180,in=-90] (d.south);
\draw[mmedged] (rev)--(note);
\end{mermaidfig}
\figcaption{Figure 21. Document quality gates. A draft that fails consistency, data
support, or formatting checks is revised rather than submitted, preventing the
amendments, regulator questions, IRB revisions, or extra review cycle that
inconsistent documents cause.}

\subsection{Prior Single-Prompt Multi-File Developments as Evidence}

This project is the latest in a line of single-prompt, multi-file builds. The Phase
2 protocol directory \cite{phase2protocol} is the most directly analogous: like this
paper, it was grown from Mermaid figures through draft, full, and final LaTeX under
one prompt. The law and bill repositories did the same for legislative text:
the review \cite{lawnarrative}, adoption \cite{adoptionguide}, and enactment
\cite{lawenactment} directories, and the two auto-bill directories
\cite{hr9510v5, bill01v4}. Table 5 lists them.

\vspace{0.2em}
\noindent\begin{tabularx}{\textwidth}{L{4.4cm} L{3.6cm} Y}
\toprule
\textbf{Repository / directory} & \textbf{Output type} & \textbf{DOI} \\
\midrule
\texttt{trial-phase-2} & Phase 2 protocol (mermaid/draft/full/final) & \href{https://doi.org/10.5281/zenodo.20807027}{10.5281/zenodo.20807027} \\
\texttt{law-physical-ai-trials/review} & Narrative case for verified trials & \href{https://doi.org/10.5281/zenodo.20685379}{10.5281/zenodo.20685379} \\
\texttt{law-physical-ai-trials/adoption} & Oncology Trial PI LLM Adoption Guide & \href{https://doi.org/10.5281/zenodo.20843290}{10.5281/zenodo.20843290} \\
\texttt{law-physical-ai-trials/enactment} & Framework for enacting H. R. 9510 & \href{https://doi.org/10.5281/zenodo.20726461}{10.5281/zenodo.20726461} \\
\texttt{single-prompt-bill/auto-bill-02} & H. R. 9510 Bill v5.0 & \href{https://doi.org/10.5281/zenodo.20619762}{10.5281/zenodo.20619762} \\
\texttt{single-prompt-bill/auto-bill-01} & H. R. 9510 Bill v4.0 & \href{https://doi.org/10.5281/zenodo.20576907}{10.5281/zenodo.20576907} \\
\bottomrule
\end{tabularx}
\figcaption{Table 5. Prior single-prompt multi-file developments cited as evidence
for the analogous approach used here.}

\subsection{2025 Application Papers Establishing LLM Trust}

LLM trust for this application was established across 2024 to 2026 (Figure 22,
Table 6). The 2025 application papers progressed from drug-cost reasoning and
clinical decision support \cite{bevacds} through glioblastoma meta-analysis
\cite{gbmmeta10yr}, lung adenocarcinoma \cite{lungadeno}, PDAC digital-twin
proposals \cite{kawchak_2025_15735068}, a 100,000-patient in silico Phase III
\cite{kawchak_2025_16415815}, QSP simulation \cite{kawchak_2025_17001137}, FDA
digital-twin compliance \cite{kawchak_2025_17239510}, end-to-end oncology-trial LLM
efficiency \cite{e2eoncoeff}, glioblastoma drug-synergy and patient-matching
machine learning \cite{gbmsynergy, aipeerreview}, and a 16-LLM code-generation
competition \cite{codegen16}.

\begin{mermaidfig}
\node[mmin] (y1) at (0,0) {2024\\LMM chem + Cancer vs CAI};
\node[mmin] (y2) at (3.5,0) {2025 H1\\cost, CDS, GBM meta, lung};
\node[mmproc] (y3) at (7.2,0) {2025 H2\\PDAC twin, 100k Phase III};
\node[mmproc] (y4) at (10.9,0) {Dec 2025\\16-LLM code competition};
\node[mmproc] (y5) at (1.8,-2.2) {Mar 2026\\National Platform};
\node[mmproc] (y6) at (5.6,-2.2) {2026\\21 CFR, H. R. 9510 v5.0};
\node[mmproc] (y7) at (9.4,-2.2) {Jun 2026\\Phase 1 + Phase 2 protocols};
\node[mmaccent] (this) at (5.6,-4.2) {This paper (v1.0)\\single-prompt document generation};
\draw[mmedge] (y1)--(y2); \draw[mmedge] (y2)--(y3); \draw[mmedge] (y3)--(y4);
\draw[mmedge] (y4) to[out=-90,in=0] (y7.east);
\draw[mmedge] (y7)--(y6); \draw[mmedge] (y6)--(y5);
\draw[mmedge] (y5) to[out=-90,in=180] (this.west);
\draw[mmedge] (y7) to[out=-90,in=0] (this.east);
\end{mermaidfig}
\figcaption{Figure 22. The author LLM-trust evidence timeline from 2024 to 2026,
terminating in this single-prompt document-generation paper.}

\vspace{0.2em}
\noindent\begin{tabularx}{\textwidth}{L{2.0cm} L{5.0cm} Y}
\toprule
\textbf{Date} & \textbf{Work} & \textbf{Contribution to trust} \\
\midrule
2025-06 & PDAC digital-twin proposals \cite{kawchak_2025_15735068} & End-to-end PDAC trial proposal scaffolding \\
2025-07 & 100,000-patient in silico Phase III \cite{kawchak_2025_16415815} & Large-scale simulated trial generation \\
2025-08 & QSP metastatic PDAC simulation \cite{kawchak_2025_17001137} & Code, VVUQ, and playbook from protocol to prediction \\
2025-10 & FDA digital-twin compliance; oncology LLM efficiency \cite{kawchak_2025_17239510, e2eoncoeff} & Compliance-aware, efficient document workflows \\
2025-11/12 & GBM drug synergy; AI peer review; 16-LLM competition \cite{gbmsynergy, aipeerreview, codegen16} & Verified, comparative LLM code generation \\
\bottomrule
\end{tabularx}
\figcaption{Table 6. Representative 2025 application papers and their contribution to
the LLM trust on which this project builds.}
